% =====================================================================
% 11. Challenges, Risks and Issues
% =====================================================================
\section{Challenges, Risks and Issues}
\label{sec:risks}

\subsection{Issues encountered}\label{subsec:issues-encountered}
\todoin{Describe real problems met so far and how they were handled, for example waiting on rig interface access, IT approval, or tooling problems. One short paragraph each.}

The main structural challenge so far has been separating the research from the company project.
The partner's plan centres on deployment, training and go-live, while the dissertation must answer four research questions.
This was resolved by separating the two sets of deliverables (Table~\ref{tab:deliverables}) and scheduling the experiments explicitly (Section~\ref{sec:timeline}).

\subsection{Risk register}\label{subsec:risk-register}
\begin{table}[H]
    \centering\small
    \caption{Risk register (L = likelihood, I = impact; H/M/L).}
    \label{tab:risks}
    \begin{tabularx}{\textwidth}{YccYY}
        \toprule
        Risk                                                          & L & I & Mitigation                                                   & Fallback                                                                           \\
        \midrule
        Access to the STT rig or data interface is delayed or refused & H & H & Adapter pattern; simulator and public data from the start & All experiments on replayed and synthetic traces; live adapter becomes future work \\
        Partner confidentiality limits what can be reported           & M & H & Agree terms early; isolate rig-specific code in one adapter & Report synthetic and public-data results only; describe the rig generically \\
        Company deliverables crowd out research work                  & H & M & Separate deliverables; protect experiment weeks in the plan & Cut demo polish before cutting experiments \\
        Analytical-only interface evaluation judged weak              & M & M & Combine KLM, clutter and accessibility; state limits plainly & Narrow RQ4 to task time and accessibility \\
        Performance targets not met                                   & M & M & Downsampling at the gateway; benchmark early                 & Report the achieved envelope honestly; reduce channel count per view               \\
        Clutter-metric or KLM tooling takes longer than planned       & M & L & Prototype on one screen early & Use one metric \\
        Scope creep from partner feature requests                     & M & M & Change requests go through the issue board with priority & Defer to post-dissertation backlog \\
        Illness or time loss                                          & L & H & Two-week buffer before submission                            & Minimum viable deliverable (below)                                                 \\
        \bottomrule
    \end{tabularx}
\end{table}

\subsection{Minimum viable deliverable}\label{subsec:minimum-viable-deliverable}
If the risks above materialise, the minimum deliverable is: the simulator and replay adapters; the processing layer with at least three filters and the alert engine; the engineer view and one other role view; and Experiments~1--3 run on synthetic and public data.
Experiment~4 would then be reduced to the accessibility audit and KLM for the engineer view.
